\documentclass[a4paper]{article}
\usepackage[affil-it]{authblk}
\usepackage[backend=bibtex,style=numeric]{biblatex}
\usepackage{geometry}
\usepackage{amsmath}
\usepackage{float, caption}
\usepackage{graphicx}
\geometry{margin=1.5cm, vmargin={0pt,1cm}}
\setlength{\topmargin}{-1cm}
\setlength{\paperheight}{29.7cm}
\setlength{\textheight}{25.3cm}

\addbibresource{citation.bib}

\begin{document}
% =================================================
\title{Numerical Analysis homework \# 6}

\author{Guang Zhang 3230105121
  \thanks{Electronic address: \texttt{zhuiguang49@foxmail.com}}}
\affil{(Information and Computing Science, Class XJ2301), Zhejiang University }

\date{Finish time: \today}

\maketitle


\section*{I. Question I}

\subsection*{I.a}
We seek to approximate the integral $I = \int_{-1}^{1} y(t) dt$ using a cubic Hermite interpolating polynomial $p_3(t)$ satisfying the conditions:
\[
p_3(-1) = y(-1), \quad p_3(0) = y(0), \quad p_3'(0) = y'(0), \quad p_3(1) = y(1).
\]
Let the polynomial be explicitly defined as $p_3(t) = a_0 + a_1 t + a_2 t^2 + a_3 t^3$.
Integrating this polynomial over $[-1, 1]$:
\[
\int_{-1}^{1} p_3(t) dt = \int_{-1}^{1} (a_0 + a_1 t + a_2 t^2 + a_3 t^3) dt.
\]
Due to the symmetry of the interval $[-1, 1]$, the integrals of the odd power terms ($t$ and $t^3$) vanish. We are left with:
\[
\int_{-1}^{1} p_3(t) dt = \left[ a_0 t + \frac{1}{3}a_2 t^3 \right]_{-1}^{1} = 2a_0 + \frac{2}{3}a_2.
\]
Now we determine the coefficients $a_0$ and $a_2$ using the interpolation conditions:
\begin{enumerate}
    \item At $t=0$: $p_3(0) = a_0 = y(0)$.
    \item At $t=1$ and $t=-1$:
    \[
    \begin{cases}
        p_3(1) = a_0 + a_1 + a_2 + a_3 = y(1) \\
        p_3(-1) = a_0 - a_1 + a_2 - a_3 = y(-1)
    \end{cases}
    \]
    Adding these two equations eliminates $a_1$ and $a_3$:
    \[
    y(1) + y(-1) = 2a_0 + 2a_2.
    \]
    Substituting $a_0 = y(0)$, we solve for $a_2$:
    \[
    2a_2 = y(1) + y(-1) - 2y(0) \implies a_2 = \frac{1}{2}[y(1) + y(-1)] - y(0).
    \]
\end{enumerate}
Substituting $a_0$ and $a_2$ back into the integral expression:
\begin{align*}
\int_{-1}^{1} p_3(t) dt &= 2y(0) + \frac{2}{3} \left( \frac{1}{2}[y(1) + y(-1)] - y(0) \right) \\
&= 2y(0) + \frac{1}{3}[y(1) + y(-1)] - \frac{2}{3}y(0) \\
&= \frac{4}{3}y(0) + \frac{1}{3}[y(1) + y(-1)] \\
&= \frac{1}{3} [y(-1) + 4y(0) + y(1)].
\end{align*}
This is exactly the formulation of Simpson's rule. Actually, we can obtain Simpson's rule by Lagrange base function, we have 
\begin{align}
  l_0(x)=\frac{x(x-1)}{2}\\
  l_1(x)=\frac-x^2+1\\
  l_2(x)=\frac{x(x+1)}{2}
\end{align}
From which we can obtain that $A_0=\frac{1}{3}, A_1=\frac{4}{3},A_2=\frac{1}{3}$, which leads to Simpson's rule:
\begin{equation}
  \int_{-1}^1 y(t)dt=\frac{1}{3}[y(-1)+4y(0)+y(14)]
\end{equation}
\subsection*{I.b}
The error of the Hermite interpolation $p_3(t)$ for a function $y \in C^4[-1, 1]$ is given by:
\[
y(t) - p_3(t) = \frac{y^{(4)}(\xi_t)}{4!} \Omega(t),
\]
where the node polynomial is $\Omega(t) = (t - (-1))(t - 0)^2 (t - 1) = (t^2 - 1)t^2$.
The error of the quadrature is obtained by integrating this remainder:
\[
E^S(y) = \int_{-1}^{1} [y(t) - p_3(t)] dt = \int_{-1}^{1} \frac{y^{(4)}(\xi_t)}{24} t^2(t^2 - 1) dt.
\]
We observe that on the interval $[-1, 1]$, the term $t^2 \ge 0$ and $t^2 - 1 \le 0$. Thus, the kernel function $K(t) = t^2(t^2 - 1)$ does not change sign (it is non-positive). By the Mean Value Theorem for Integrals, there exists some $\zeta \in (-1, 1)$ such that:
\[
E^S(y) = \frac{y^{(4)}(\zeta)}{24} \int_{-1}^{1} (t^4 - t^2) dt.
\]
Computing the definite integral:
\[
\int_{-1}^{1} (t^4 - t^2) dt = \left[ \frac{t^5}{5} - \frac{t^3}{3} \right]_{-1}^{1} = 2\left(\frac{1}{5} - \frac{1}{3}\right) = 2\left(\frac{3-5}{15}\right) = -\frac{4}{15}.
\]
Therefore, the error is:
\[
E^S(y) = \frac{y^{(4)}(\zeta)}{24} \left( -\frac{4}{15} \right) = -\frac{1}{90} y^{(4)}(\zeta).
\]

\subsection*{I.c}
To derive the composite Simpson's rule for $\int_{a}^{b} f(x) dx$, we partition $[a, b]$ into $n$ subintervals (where $n$ is even) with step size $h = \frac{b-a}{n}$ and nodes $x_k = a + kh$. We consider the sub-interval $[x_{2k}, x_{2k+2}]$ of width $2h$.

We perform a change of variables $x = x_{2k+1} + th$, where $t \in [-1, 1]$. Then $dx = h dt$. Let $y(t) = f(x_{2k+1} + th)$.
The integral over one panel is:
\[
I_k = \int_{x_{2k}}^{x_{2k+2}} f(x) dx = h \int_{-1}^{1} y(t) dt.
\]
Using the result from (a), the approximation is:
\[
I_k \approx h \cdot \frac{1}{3} [y(-1) + 4y(0) + y(1)] = \frac{h}{3} [f(x_{2k}) + 4f(x_{2k+1}) + f(x_{2k+2})].
\]
Summing over all $k = 0, \dots, n/2 - 1$, we get the composite formula:
\[
I_n^S(f) = \sum_{k=0}^{n/2-1} I_k = \frac{h}{3} \left[ f(a) + 4\sum_{k=0}^{n/2-1} f(x_{2k+1}) + 2\sum_{k=1}^{n/2-1} f(x_{2k}) + f(b) \right].
\]

\textbf{Error Estimate:}
From part (b), the error on the standard interval is $E^S(y) = -\frac{1}{90} y^{(4)}(\zeta_k)$.
By the Chain Rule, derivatives transform as $y^{(4)}(t) = h^4 f^{(4)}(x)$.
Thus, the error for the $k$-th sub-interval is:
\[
E_k = h \cdot E^S(y) = h \left( -\frac{1}{90} h^4 f^{(4)}(\xi_k) \right) = -\frac{h^5}{90} f^{(4)}(\xi_k), \quad \xi_k \in (x_{2k}, x_{2k+2}).
\]
The total error is the sum of errors over the $n/2$ sub-intervals:
\[
E_{total} = \sum_{k=0}^{n/2-1} E_k = -\frac{h^5}{90} \sum_{k=0}^{n/2-1} f^{(4)}(\xi_k).
\]
Using the Intermediate Value Theorem, since $f^{(4)}$ is continuous, $\sum_{k=0}^{n/2-1} f^{(4)}(\xi_k) = \frac{n}{2} f^{(4)}(\xi)$ for some $\xi \in (a, b)$.
Substituting $n = \frac{b-a}{h}$:
\[
E_{total} = -\frac{h^5}{90} \cdot \frac{n}{2} f^{(4)}(\xi) = -\frac{h^5}{90} \cdot \frac{b-a}{2h} f^{(4)}(\xi) = -\frac{b-a}{180} h^4 f^{(4)}(\xi).
\]


\section*{II. Question II}
We aim to approximate the integral $I = \int_{0}^{1} e^{-x^2} dx$ with an absolute error strictly less than $\epsilon = 0.5 \times 10^{-6}$.
Let $f(x) = e^{-x^2}$. The interval is $[a, b] = [0, 1]$.

\subsection*{II.a}
The error bound for the composite trapezoidal rule is given by:
\[
|E_n^T| \le \frac{(b-a)^3}{12n^2} \max_{x \in [a, b]} |f''(x)|.
\]
First, we compute the second derivative of $f(x)$:
\begin{align*}
f'(x) &= -2x e^{-x^2} \\
f''(x) &= (4x^2 - 2)e^{-x^2}
\end{align*}
By caculation $M_2 = \max_{x \in [0, 1]} |f''(x)|=2$.
Substituting into the error inequality:
\begin{align*}
\frac{(1-0)^3}{12n^2} \cdot 2 &< 0.5 \times 10^{-6} \\
\frac{1}{6n^2} &< 5 \times 10^{-7} \\
n^2 &> \frac{1}{30 \times 10^{-7}} \approx 333,333.33 \\
n &> \sqrt{333,333.33} \approx 577.35
\end{align*}
Since $n$ must be an integer, the required number of subintervals is:
\[
n = 578.
\]

\subsection*{II.b}
The error bound for the composite Simpson's rule is given by:
\[
|E_n^S| \le \frac{(b-a)^5}{180n^4} \max_{x \in [a, b]} |f^{(4)}(x)|.
\]
We determine the fourth derivative $f^{(4)}(x)$:
\begin{align*}
f'''(x) &= (8x - 8x^3 + 4x)e^{-x^2} = (12x - 8x^3)e^{-x^2} \\
f^{(4)}(x) &= 4(4x^4 - 12x^2 + 3)e^{-x^2}
\end{align*}
By caculation $M_4 = \max_{x \in [0, 1]} |f^{(4)}(x)|=12$.
Substituting into the error inequality:
\begin{align*}
\frac{(1-0)^5}{180n^4} \cdot 12 &< 0.5 \times 10^{-6} \\
\frac{1}{15n^4} &< 5 \times 10^{-7} \\
n^4 &> \frac{1}{75 \times 10^{-7}} \approx 133,333.33 \\
n &> \sqrt[4]{133,333.33} \approx 19.1
\end{align*}
For Simpson's rule, $n$ must be an even integer. The smallest even integer greater than $19.1$ is:
\[
n = 20.
\]


\section*{III. Question III}

\subsection*{III.a}

We seek a polynomial $\pi_2(t) = t^2 + at + b$ that is orthogonal to $P_1$ with respect to the weight function $\rho(t) = e^{-t}$ on $[0, +\infty)$.
The orthogonality conditions are:
\[
\langle \pi_2, 1 \rangle = 0 \quad \text{and} \quad \langle \pi_2, t \rangle = 0.
\]
Using the property $\int_{0}^{+\infty} t^m e^{-t} dt = m!$, we compute the inner products:
\begin{align}
  \int_{0}^{+\infty} (t^2 + at + b) e^{-t} dt = 2! + a(1!) + b(0!) = 2 + a + b = 0\\
  \int_{0}^{+\infty} t(t^2 + at + b) e^{-t} dt = \int_{0}^{+\infty} (t^3 + at^2 + bt) e^{-t} dt = 3! + 2!a + 1!b = 6 + 2a + b = 0
\end{align}
Then we solve the linear system which leads to $a=-4,b=2$.
Thus, the orthogonal polynomial is:
\[
\pi_2(t) = t^2 - 4t + 2.
\]

\subsection*{III.b }
\textbf{1. Finding the Nodes:}
The nodes $t_1, t_2$ are the roots of $\pi_2(t) = 0$:
\[
t = \frac{4 \pm \sqrt{16 - 8}}{2} = 2 \pm \sqrt{2}.
\]
Let $t_1 = 2 - \sqrt{2}$ and $t_2 = 2 + \sqrt{2}$.\newline
\textbf{2. Finding the Weights:}
The quadrature formula is exact for degree $2n-1 = 3$. We can determine the weights $w_1, w_2$ by enforcing exactness for $f(t) = 1$ and $f(t) = t$:
\[
\begin{cases}
w_1 + w_2 = \int_{0}^{+\infty} e^{-t} dt = 1 \\
w_1 t_1 + w_2 t_2 = \int_{0}^{+\infty} t e^{-t} dt = 1
\end{cases}
\]
Solving this system:
\[
w_1 (2-\sqrt{2}) + w_2 (2+\sqrt{2}) = 1 \implies 2(w_1+w_2) + \sqrt{2}(w_2-w_1) = 1.
\]
Since $w_1+w_2=1$, we have $2 + \sqrt{2}(w_2-w_1) = 1 \implies w_2 - w_1 = -\frac{1}{\sqrt{2}}$.
Combined with $w_1+w_2=1$, we get:
\[
w_1 = \frac{2+\sqrt{2}}{4}, \quad w_2 = \frac{2-\sqrt{2}}{4}.
\]
The formula is:
\[
\int_{0}^{+\infty} f(t)e^{-t} dt \approx \frac{2+\sqrt{2}}{4} f(2-\sqrt{2}) + \frac{2-\sqrt{2}}{4} f(2+\sqrt{2}).
\]
\textbf{3. Error Term:}
The remainder for Gaussian quadrature is given by:
\[
E_2(f) = \frac{f^{(4)}(\tau)}{(2n)!} \int_{0}^{+\infty} [\pi_2(t)]^2 e^{-t} dt.
\]
First, compute the integral constant:
\[
\int_{0}^{+\infty} (t^2 - 4t + 2)^2 e^{-t} dt = \int_{0}^{+\infty} (t^4 - 8t^3 + 20t^2 - 16t + 4) e^{-t} dt.
\]
Using $\int t^m e^{-t} = m!$:
\[
= 24 - 8(6) + 20(2) - 16(1) + 4 = 24 - 48 + 40 - 16 + 4 = 4.
\]
Thus, the error term is:
\[
E_2(f) = \frac{f^{(4)}(\tau)}{24} \cdot 4 = \frac{1}{6} f^{(4)}(\tau), \quad \tau > 0.
\]

\subsection*{III.c}
We approximate $I = \int_{0}^{+\infty} \frac{1}{1+t} e^{-t} dt$. Here $f(t) = \frac{1}{1+t}$.
\textbf{1. Approximation:}
Using the formula derived in (b):
\begin{align*}
I_{approximate} &= w_1 f(t_1) + w_2 f(t_2) \\
&= \frac{2+\sqrt{2}}{4} \cdot \frac{1}{1 + (2-\sqrt{2})} + \frac{2-\sqrt{2}}{4} \cdot \frac{1}{1 + (2+\sqrt{2})} \\
&= \frac{2+\sqrt{2}}{4(3-\sqrt{2})} + \frac{2-\sqrt{2}}{4(3+\sqrt{2})}
\end{align*}
Simplifying the terms:
\begin{align*}
\frac{2+\sqrt{2}}{4(3-\sqrt{2})} \cdot \frac{3+\sqrt{2}}{3+\sqrt{2}} &= \frac{6 + 5\sqrt{2} + 2}{4(9-2)} = \frac{8+5\sqrt{2}}{28} \\
\frac{2-\sqrt{2}}{4(3+\sqrt{2})} \cdot \frac{3-\sqrt{2}}{3-\sqrt{2}} &= \frac{6 - 5\sqrt{2} + 2}{4(9-2)} = \frac{8-5\sqrt{2}}{28}
\end{align*}
Summing them up:
\[
I_{approx} = \frac{16}{28} = \frac{4}{7} \approx 0.57142857.
\]
\textbf{2. Comparison and Identification of $\tau$:}
Given the exact value $I_{true} = 0.59634736$, the true error is:
\[
E = I_{true} - I_{approx} = 0.59634736 - 0.57142857 = 0.02491879.
\]
From (b), the theoretical error is $E_2(f) = \frac{1}{6} f^{(4)}(\tau)$.
We compute derivatives of $f(t) = (1+t)^{-1}$:
\[
f'(t) = -(1+t)^{-2}, \quad \dots, \quad f^{(4)}(t) = 24(1+t)^{-5}.
\]
Thus, $E_2(f) = \frac{1}{6} \cdot 24(1+\tau)^{-5} = 4(1+\tau)^{-5}$.
Equating the true error to the theoretical form:
\[
4(1+\tau)^{-5} = 0.02491879 \implies (1+\tau)^5 = \frac{4}{0.02491879} \approx 160.521.
\]
Solving for $\tau$:
\[
1+\tau = (160.521)^{1/5} \approx 2.766 \implies \tau \approx 1.76125
\]
This value $\tau \approx 1.76125$ lies within the integration domain $(0, +\infty)$, confirming the validity of the error term.



\section*{IV. Question IV}

\subsection*{IV.a}
We seek basis functions of the form:
\[
h_m(t) = (a_m + b_m t) l_m^2(t), \quad q_m(t) = (c_m + d_m t) l_m^2(t).
\]
Recall the properties of the Lagrange basis $l_m(t)$:
\[
l_m(x_j) = \delta_{mj}.
\]
Differentiation yields $h_m'(t) = b_m l_m^2(t) + (a_m + b_m t) 2 l_m(t) l_m'(t)$.
\textbf{1. Determining $h_m(t)$:}
We require $h_m(x_m) = 1$ and $h_m'(x_m) = 0$.
\begin{itemize}
    \item $h_m(x_m) = (a_m + b_m x_m) \cdot 1^2 = 1 \implies a_m + b_m x_m = 1$.
    \item $h_m'(x_m) = b_m \cdot 1 + (a_m + b_m x_m) \cdot 2 l_m'(x_m) = 0$.
\end{itemize}
Substituting $a_m + b_m x_m = 1$ into the derivative equation:
\[
b_m + 2(1) l_m'(x_m) = 0 \implies b_m = -2 l_m'(x_m).
\]
Then $a_m = 1 - b_m x_m = 1 + 2 x_m l_m'(x_m)$.
Thus, the coefficients are $a_m = 1 + 2 x_m l_m'(x_m)$ and $b_m = -2 l_m'(x_m)$.
Alternatively written:
\[
h_m(t) = [1 - 2(t - x_m) l_m'(x_m)] l_m^2(t).
\]
\textbf{2. Determining $q_m(t)$:}
We require $q_m(x_m) = 0$ and $q_m'(x_m) = 1$.
\begin{itemize}
    \item $q_m(x_m) = (c_m + d_m x_m) \cdot 1 = 0 \implies c_m = -d_m x_m$.
    \item $q_m'(x_m) = d_m + 2(c_m + d_m x_m) l_m'(x_m) = 1$.
\end{itemize}
Since $c_m + d_m x_m = 0$, the second term vanishes:
\[
d_m + 0 = 1 \implies d_m = 1.
\]
Then $c_m = -x_m$.
Thus, the coefficients are $c_m = -x_m$ and $d_m = 1$.
This gives:
\[
q_m(t) = (t - x_m) l_m^2(t).
\]

\subsection*{IV.b}
Let $p(t)$ be the Hermite interpolating polynomial for $f(t)$. Integrating $p(t)$ gives the quadrature rule:
\[
I_n(f) = \int_a^b p(t) \rho(t) dt = \sum_{k=1}^n \left[ f(x_k) \int_a^b h_k(t) \rho(t) dt + f'(x_k) \int_a^b q_k(t) \rho(t) dt \right].
\]
Defining the weights and coefficients as:
\[
w_k = \int_a^b h_k(t) \rho(t) dt, \quad \mu_k = \int_a^b q_k(t) \rho(t) dt.
\]
We obtain the rule:
\[
I_n(f) = \sum_{k=1}^n [w_k f(x_k) + \mu_k f'(x_k)].
\]
Since $p(t)$ reproduces any polynomial in ${P}_{2n-1}$ exactly, the error $E_n(p) = 0$ for all $p \in {P}_{2n-1}$.

\subsection*{IV.c}
To eliminate the derivative terms (i.e., $\mu_k = 0$ for all $k$), we examine the expression for $\mu_k$:
\[
\mu_k = \int_a^b (t - x_k) l_k^2(t) \rho(t) dt.
\]
Let $v_n(t) = \prod_{j=1}^n (t - x_j)$ be the node polynomial. Recall that the Lagrange basis can be written as $l_k(t) = \frac{v_n(t)}{(t - x_k) v_n'(x_k)}$.
Substituting one factor of $l_k(t)$ into the integral:
\[
(t - x_k) l_k^2(t) = (t - x_k) \frac{v_n(t)}{(t - x_k) v_n'(x_k)} l_k(t) = \frac{1}{v_n'(x_k)} v_n(t) l_k(t).
\]
Thus,
\[
\mu_k = \frac{1}{v_n'(x_k)} \int_a^b v_n(t) l_k(t) \rho(t) dt.
\]
For $\mu_k$ to vanish for all $k$, the integral $\int_a^b v_n(t) l_k(t) \rho(t) dt$ must be zero.
Since $\{l_k(t)\}_{k=1}^n$ forms a basis for ${P}_{n-1}$, the condition $\int v_n l_k \rho dt = 0$ for all $k$ is equivalent to saying that $v_n(t)$ is orthogonal to the entire space ${P}_{n-1}$ with respect to the weight $\rho(t)$.

\textbf{Conclusion:}
The condition required is that the node polynomial $v_n(t)$ must be the \textbf{orthogonal polynomial} of degree $n$ with respect to the weight function $\rho(t)$. This implies that the nodes $x_k$ must be the roots of this orthogonal polynomial (i.e., they must be the \textbf{Gauss nodes}). Under this condition, the quadrature formula becomes the standard Gaussian Quadrature with degree of exactness $2n-1$.

\section*{V.Question V}

\subsection*{V.1 Proof of Lemma 6.40}
We consider the second-order central difference formula for the second derivative:
\[
D^2u(\bar{x}) = \frac{u(\bar{x}-h) - 2u(\bar{x}) + u(\bar{x}+h)}{h^2}.
\]
\textbf{1. Truncation Error:}
Assuming $u \in C^4$, we expand $u(\bar{x} \pm h)$ using Taylor series about $\bar{x}$:
\[
u(\bar{x} \pm h) = u(\bar{x}) \pm h u'(\bar{x}) + \frac{h^2}{2} u''(\bar{x}) \pm \frac{h^3}{6} u'''(\bar{x}) + \frac{h^4}{24} u^{(4)}(\xi_\pm).
\]
Summing the two expansions:
\[
u(\bar{x}+h) + u(\bar{x}-h) = 2u(\bar{x}) + h^2 u''(\bar{x}) + \frac{h^4}{24} (u^{(4)}(\xi_+) + u^{(4)}(\xi_-)).
\]
By the Intermediate Value Theorem, $\frac{1}{2}(u^{(4)}(\xi_+) + u^{(4)}(\xi_-)) = u^{(4)}(\xi)$ for some $\xi \in (\bar{x}-h, \bar{x}+h)$. Thus:
\[
u(\bar{x}+h) - 2u(\bar{x}) + u(\bar{x}-h) = h^2 u''(\bar{x}) + \frac{h^4}{12} u^{(4)}(\xi).
\]
Dividing by $h^2$ gives the truncation error:
\[
D^2u(\bar{x}) - u''(\bar{x}) = \frac{h^2}{12} u^{(4)}(\xi).
\]
\textbf{2. Rounding Error:}
Let $\tilde{u}$ be the computed values with rounding error $|\tilde{u} - u| \le \epsilon$. The computed difference is:
\[
D^2\tilde{u} = \frac{\tilde{u}(\bar{x}-h) - 2\tilde{u}(\bar{x}) + \tilde{u}(\bar{x}+h)}{h^2}.
\]
The error due to rounding is bounded by:
\[
|E_{round}| \le \frac{|\epsilon| + |-2\epsilon| + |\epsilon|}{h^2} = \frac{4\epsilon}{h^2}.
\]
\textbf{3. Total Error Bound:}
Combining both errors, the total error bound is:
\[
|E(h)| \le \frac{h^2}{12} M_4 + \frac{4\epsilon}{h^2}, \quad \text{where } M_4 = \max |u^{(4)}(x)|.
\]
\textbf{4. Minimizing the Error:}
To find the optimal $h$, we define $\phi(h) = \frac{M_4}{12} h^2 + \frac{4\epsilon}{h^2}$ and set $\phi'(h) = 0$:
\[
\frac{M_4 h}{6} - \frac{8\epsilon}{h^3} = 0 \implies h^4 = \frac{48\epsilon}{M_4} \implies h_{opt} = \left( \frac{48\epsilon}{M_4} \right)^{1/4}.
\]
The minimal error scales as $O(\epsilon^{1/2})$.

\subsection*{V.2 Designing a Fourth-Order Symmetric Formula}
We construct a formula using a 5-point symmetric stencil: $\bar{x}, \bar{x} \pm h, \bar{x} \pm 2h$.
\[
D_{4}^2 u(\bar{x}) = \frac{c_0 u(\bar{x}) + c_1 [u(\bar{x}+h) + u(\bar{x}-h)] + c_2 [u(\bar{x}+2h) + u(\bar{x}-2h)]}{h^2}.
\]
Expanding the terms using Taylor series:
\begin{align*}
u(\bar{x} \pm h) + u(\bar{x} \mp h) &= 2u + h^2 u'' + \frac{h^4}{12} u^{(4)} + \frac{h^6}{360} u^{(6)} + \dots \\
u(\bar{x} \pm 2h) + u(\bar{x} \mp 2h) &= 2u + 4h^2 u'' + \frac{16h^4}{12} u^{(4)} + \frac{64h^6}{360} u^{(6)} + \dots
\end{align*}
Substitute into the formula and match coefficients to approximate $1 \cdot u''$:
\begin{align}
u: & c_0 + 2c_1 + 2c_2 = 0 \\
u'': & c_1 + 4c_2 = 1 \\
u^{(4)}: & \frac{1}{12} c_1 + \frac{16}{12} c_2 = 0 \implies c_1 + 16c_2 = 0
\end{align}
We can finally obtain that the formula is:
\[
D_4^2 u(\bar{x}) = \frac{-\frac{1}{12}u_{\pm 2} + \frac{16}{12}u_{\pm 1} - \frac{30}{12}u_0}{h^2} = \frac{-u_{\pm 2} + 16u_{\pm 1} - 30u_0}{12h^2}.
\]
\newline
\textbf{1. Truncation Error Bound:}
The leading error term comes from $h^6$:
\[
\text{Term } = \frac{1}{h^2} \left[ c_1 \frac{h^6}{360} + c_2 \frac{64h^6}{360} \right] u^{(6)} = \frac{h^4}{360} \left[ \frac{16}{12} + 64(-\frac{1}{12}) \right] u^{(6)}
\]
\[
= \frac{h^4}{360} \left[ \frac{16 - 64}{12} \right] u^{(6)} = \frac{h^4}{360} (-4) u^{(6)} = -\frac{1}{90} h^4 u^{(6)}.
\]
Bound: $\frac{M_6}{90} h^4$.\newline
\textbf{2. Rounding Error Bound:}
Sum of absolute coefficients in numerator: $|-1| + |16| + |-30| + |16| + |-1| = 64$.
Denominator is $12h^2$.
\[
|E_{round}| \le \frac{64\epsilon}{12h^2} = \frac{16\epsilon}{3h^2}.
\]
\textbf{3. Minimizing the Total Error:}
\[
\psi(h) = \frac{M_6}{90} h^4 + \frac{16\epsilon}{3h^2}.
\]
Set $\psi'(h) = 0$:
\[
\frac{4 M_6 h^3}{90} - \frac{32\epsilon}{3h^3} = 0 \implies \frac{2 M_6 h^6}{45} = \frac{32\epsilon}{3}
\]
\[
h^6 = \frac{32 \cdot 45 \epsilon}{6 M_6} = \frac{240\epsilon}{M_6} \implies h_{opt} \approx \left( \frac{240\epsilon}{M_6} \right)^{1/6}.
\]

\subsection*{V.3 Comparison and Observation}
\begin{itemize}
    \item \textbf{Second-Order Case:} Optimal step size $h \sim \epsilon^{1/4}$, minimal error $E \sim \epsilon^{1/2}$.
    \item \textbf{Fourth-Order Case:} Optimal step size $h \sim \epsilon^{1/6}$, minimal error $E \sim \epsilon^{2/3}$.
\end{itemize}
\textbf{Observation:} Since machine epsilon $\epsilon$ is very small (typically $10^{-16}$), $\epsilon^{2/3}$ is significantly smaller than $\epsilon^{1/2}$. This means the fourth-order method achieves much higher accuracy. Additionally, since $\epsilon^{1/6} > \epsilon^{1/4}$ for small $\epsilon$, the optimal step size for the higher-order method is \textbf{larger}, meaning we can achieve better accuracy without taking an extremely small step size (which avoids catastrophic cancellation issues).
\newline
\newline\newline\newline
\textbf{Note:} As the final examinations approach, I am busy preparing for the exams, so the latex code is actually given by Gemini3, while the solution is all done by hand myself in advance. 

\end{document}